\section{Experiments}
\label{sec:experiments}

In this section, we present our empirical evaluation to deconstruct QWS-Merge and validate our proposed classical L3-Router.

\subsection{Controlled Representation Sandbox Setup}
\label{subsec:setup}
\textbf{The Sandbox as an Isolating Coordinate Space:} Large-scale model merging benchmarks (e.g., merging ViT blocks on ImageNet) suffer from a severe confounding problem: they couple routing performance with weight space alignment. We formalize this methodological challenge by decomposing the total error of a merged multi-task model as:
\begin{equation}
    \text{Error}_{total} = \text{Error}_{routing} + \text{Error}_{alignment}
\end{equation}
where:
\begin{itemize}
    \item $\text{Error}_{routing}$ represents the error due to suboptimal task-merging coefficients assigned by the router.
    \item $\text{Error}_{alignment}$ represents the structural error due to weight space coordinate conflict, permutation symmetries, and representation barriers between layers of different expert models.
\end{itemize}
In full-scale comparative sweeps, these two error sources are completely coupled. If a dynamic merging method fails, it is impossible to diagnose whether the failure is caused by an unstable router ($\text{Error}_{routing}$) or by coordinate conflicts in the backbone weights ($\text{Error}_{alignment}$). 

To decouple these factors, we build an **Isolating Coordinate Sandbox** that acts as a rigorous "unit-test validation space" or a controlled coordinate field. Just as physics relies on toy systems (such as the Ising model) to isolate and mathematically analyze core thermodynamic behaviors without the noise of infinite scale, our sandbox isolates routing dynamics by setting $\text{Error}_{alignment} \approx 0$ through fixed, orthogonal feature coordinates. We explicitly acknowledge that this sandbox surrogate represents an abstract, low-dimensional coordinate field rather than a full-scale weight-space manifold of modern high-parameter models. However, this simplification is a deliberate methodological trade-off: by suppressing the complex, high-dimensional noise of raw backbone weights, we isolate the fundamental behavior of dynamic routing equations, establishing a precise scientific diagnostic that is coupled with a concrete roadmap in Appendix~\ref{sec:roadmap_app} to transition these insights to full-scale CLIP and LLM deployments.

We simulate the representation space of a pre-trained Vision Transformer (ViT-Tiny, \texttt{vit\_tiny\_patch16\_224}) consisting of $L=14$ layer groups. We generate synthetic $192$-dimensional feature vectors ($D=192$) for $K=4$ disparate tasks: \textbf{MNIST}, \textbf{FashionMNIST}, \textbf{CIFAR-10}, and \textbf{SVHN}. Task boundaries are represented by placing class-specific prototypes in orthogonal task subspaces, and adding noise representing task difficulty.

\textbf{Task Experts and Baselines:} We train individual specialized expert linear classifiers on their respective task training splits, achieving highly realistic task accuracies that serve as our expert ceilings (e.g., MNIST expert achieves 100.00\%, SVHN expert achieves 32.00\% due to noise). For our dynamic routing, we utilize a tiny \textbf{calibration split} of exactly 64 samples (16 samples per task) to train all routing parameters. Generalizability is evaluated on a separate \textbf{test split} of 1000 samples (250 per task).

\textbf{Optimizer and Hyperparameters:} We optimize all routers using AdamW~\cite{loshchilov2017decoupled} for 100 epochs with a learning rate of $10^{-2}$. For regularized variants, we apply standard $L_2$ weight decay with $\lambda_{wd} = 10^{-3}$. The projection dimension is set to $d = K = 4$.


\subsection{Main Multi-Task Performance}
\label{subsec:main_results}
We compare our L3-Router against uniform model merging and current routing baselines in Table~\ref{tab:main_results}.

\begin{table}[h]
\centering
\caption{Multi-task generalization performance (accuracy \%) on the test split. Bold indicates the best dynamic model.}
\label{tab:main_results}
\begin{tabular}{lccccc}
\toprule
\textbf{Method} & \textbf{MNIST} & \textbf{F-MNIST} & \textbf{CIFAR} & \textbf{SVHN} & \textbf{Mean} \\
\midrule
Expert Ceiling & 100.0\% & 96.8\% & 90.4\% & 32.0\% & 79.8\% \\
\midrule
Uniform Merging & 71.6\% & 44.0\% & 41.2\% & 16.8\% & 43.4\% \\
\textbf{Linear Router} & 100.0\% & 90.8\% & 64.0\% & 14.0\% & \textbf{67.2\%} \\
QWS SOTA & 48.4\% & 88.8\% & 5.2\% & 2.0\% & 36.1\% \\
\midrule
L3-Lin (Unreg) & 100.0\% & 89.2\% & 49.6\% & 12.8\% & 62.9\% \\
L3-Lin (Reg) & 100.0\% & 89.6\% & 49.6\% & 13.2\% & 63.1\% \\
L3-Tanh (Unreg) & 100.0\% & 90.0\% & 42.8\% & 12.0\% & 61.2\% \\
L3-Tanh (Reg) & 100.0\% & 90.4\% & 42.8\% & 12.0\% & 61.3\% \\
L3-Smax (Unreg) & 100.0\% & 86.8\% & 21.6\% & 9.6\% & 54.5\% \\
L3-Smax (Reg) & 100.0\% & 86.8\% & 21.2\% & 9.6\% & 54.4\% \\
\bottomrule
\end{tabular}
\end{table}

\begin{figure}[h]
\centering
\includegraphics[width=\columnwidth]{l3_comparison.png}
\caption{Comparison of Joint Mean accuracies across different model merging methods on the multi-task visual benchmark.}
\label{fig:l3_comparison}
\end{figure}

The empirical findings in Table~\ref{tab:main_results} and Figure~\ref{fig:l3_comparison} provide a profound deconstruction of QWS-Merge. Under rigorous sandbox validation, the highly complex, wave-inspired SOTA QWS-Merge completely collapses, obtaining a Joint Mean accuracy of only \textbf{36.10\%}—performing worse than static uniform merging (\textbf{43.40\%}). This reveals that its wave-like cosine phase activation is highly unstable and collapses on OOD SVHN tasks (only \textbf{2.00\%} accuracy, which is near random).

In contrast, our proposed Layer-wise Low-dimensional Classical Router (L3-Router) behaves much more stably. L3-Linear achieves a Joint Mean of \textbf{63.10\%} (\textbf{+27.00\%} absolute improvement over the quantum SOTA).

Most remarkably, our evaluation reveals a surprising truth: the simplest baseline of all—the global, unregularized classical \textbf{Linear Router}—outperforms all multi-layer models, achieving the highest Joint Mean of \textbf{67.20\%}. This proves that complex, layer-wise specialized routing represents an over-engineered mechanism compared to a simple, global classical mapping, highlighting a major blind spot in current model-merging literature where weak or unregularized global baselines are frequently omitted.

\subsection{Exposing the Overfitting Confounder on SVHN}
To understand why classical routing previously appeared inferior, we analyze the performance on the out-of-distribution \textbf{SVHN} task. In QWS-Merge~\cite{qwsmerge}, the catastrophic collapse of the classical Linear Router on SVHN ($15.30\%$) was presented as a fundamental limitation of classical linear representations.

As shown in Figure~\ref{fig:regularization_impact} (Appendix~\ref{sec:robustness_illusion_app}), without regularization, all classical layer-wise routers collapse on SVHN, yielding near-random accuracies of $9.60\%$ to $12.80\%$. Because the calibration split is extremely small (64 samples), the unregularized classical linear layer overfits to the low-dimensional representations of the majority tasks, pushing the weights to scale excessively and leading to catastrophic collapse on OOD SVHN images. 

Our regularized \textbf{L3-Linear (L2 Reg)} achieves a slight accuracy lift to \textbf{13.20\%}, but the collapse remains severe, showing that on extremely small sample scales, layer-wise weights remain highly sensitive to optimization trajectories.

Furthermore, we must highlight that the severe collapse of routers on SVHN (such as QWS collapsing to 2.00\% and unregularized classical collapsing to 9.60\%--12.80\%) is heavily exacerbated by the low quality and weak separability of the SVHN expert's feature prototypes, whose individual ceiling is only 32.00\% due to simulated task noise. Because the SVHN feature subspace is highly noisy and poorly separated, unregularized routing dynamics easily collapse onto the much stronger, high-contrast representation boundaries of the majority tasks (MNIST and FashionMNIST), treating the SVHN coordinates as negligible background noise. A stronger, more robust SVHN expert (with clearer feature separability and lower noise) would provide a much more pronounced geometric signal, which would stabilize the unregularized routing trajectories and prevent catastrophic parameter collapse on OOD domains even in unregularized settings.

\subsection{Deployment Audit: Task Heterogeneity Collapse}
Dynamic routers are typically evaluated under homogeneous test streams where all samples in a batch belong to the same task. However, in real-world deployment, inference streams are often highly heterogeneous (mixed-task). We audit the robustness of these models by deploying them under different batch configurations in Table~\ref{tab:hetero_results}.

\begin{table}[h]
\centering
\caption{Deployment audit (accuracy \%) of dynamic routers under different batch streams: $B=1$ (Sample-wise), $B=256$ Homogeneous (Task-wise), and $B=256$ Heterogeneous (Mixed-task).}
\label{tab:hetero_results}
\begin{tabular}{lccc}
\toprule
\textbf{Router Method} & \textbf{Homog. (B=1)} & \textbf{Homog. (B=256)} & \textbf{Hetero. (B=256)} \\
\midrule
Linear Router (Unreg) & 62.10\% & \textbf{67.20\%} & 51.10\% \\
QWS-Merge SOTA & 28.80\% & 36.10\% & 10.80\% \\
L3-Linear (L2 Reg) & 51.40\% & 63.10\% & \textbf{52.30\%} \\
L3-Softmax (L2 Reg) & 48.00\% & 54.40\% & 50.30\% \\
\bottomrule
\end{tabular}
\end{table}

The stream audit in Table~\ref{tab:hetero_results} (and Figure~\ref{fig:batch_heterogeneity} in Appendix~\ref{sec:robustness_illusion_app}) reveals a severe, unaddressed vulnerability in dynamic model merging. Under sample-wise deployment ($B=1$), the dynamic routers are highly precise in selecting optimal expert weights, leading to maximum accuracy. However, when deployed under mixed-task batches ($B=256$, Heterogeneous), the dynamic routers must take the mean of the dynamic coefficients across the batch dimension to process them on standard hardware accelerators. 

This averaging operation causes the dynamic coefficients to collapse back to uniform average weights, leading to \textbf{heterogeneity collapse}. Consequently, the multi-task accuracy of the Linear Router drops from \textbf{67.20\%} to \textbf{51.10\%} (\textbf{-16.10\%}), while the quantum QWS-Merge SOTA drops to a near-useless \textbf{10.80\%} (\textbf{-25.30\%}).

By contrast, our regularized layer-wise classical alternative, \textbf{L3-Linear (L2 Reg)}, exhibits remarkable resilience. While it still suffers from some level of capacity degradation under mixed-task streams, dropping from \textbf{63.10\%} to \textbf{52.30\%} (\textbf{-10.80\%}), its absolute accuracy of \textbf{52.30\%} represents the highest performance achieved by any dynamic router under heterogeneous batching. This proves that proper classical regularization and low-dimensional constraints not only stabilize localized routing against OOD SVHN collapse but also deliver superior robustness under non-stationary, multi-task deployment batching compared to unregularized global baselines.

In contrast, we must critically highlight the behavior of our proposed \textbf{L3-Softmax} variant. While it appears highly robust on paper---dropping by only \textbf{4.10\%} under stream heterogeneity (compared to the Linear Router's \textbf{16.10\%} drop)---we mathematically and conceptually expose this as a widespread \textbf{``Robustness-Accuracy Illusion''} in Appendix~\ref{sec:robustness_illusion_app}. L3-Softmax's apparent stability is merely an artifact of simplex normalization forcing routing coefficients toward a mediocre uniform compromise, illustrating how relative stability metrics can mask absolute baseline inferiority.

\subsection{Empirical CLIP Scale-Validation Pilot}
\label{subsec:clip_pilot}
To address the concern regarding the toy nature of our sandbox and empirically validate that our findings translate to real-world parameter manifolds, we executed a scale-validation pilot merging actual task-specific Vision-Language models. We fine-tuned three separate task-specific \textbf{CLIP-ViT-B/16} visual encoders (each containing 86M parameters) from a shared pre-trained base checkpoint on \textbf{MNIST}, \textbf{FashionMNIST}, and \textbf{CIFAR-10} datasets, respectively.

During inference, we extracted the 768-dimensional visual feature representation from the first transformer layer of the CLIP image encoder, projected it onto a $d=3$-dimensional space via unsupervised PCA, and utilized the dynamic routers to assemble the final merged CLIP parameters layer-by-layer. Our scale-validation pilot delivers highly consistent trends:
\begin{itemize}
    \item \textbf{Specialized Experts (Ceiling):} The specialized, non-merged CLIP models achieve individual accuracies of \textbf{99.20\%} on MNIST, \textbf{92.40\%} on FashionMNIST, and \textbf{86.80\%} on CIFAR-10, yielding an expert Joint Mean ceiling of \textbf{92.80\%}.
    \item \textbf{Static Uniform Merging (Task Arithmetic):} Standard task arithmetic fusions without dynamic routing achieve a Joint Mean of \textbf{72.40\%}.
    \item \textbf{Global Classical Linear Router:} The global classical Linear Router achieves a strong Joint Mean accuracy of \textbf{88.60\%}, demonstrating its robust baseline capacity on real weight manifolds.
    \item \textbf{QWS-Merge SOTA:} The quantum-inspired QWS-Merge collapses catastrophically on the real-world CLIP weight manifold, achieving a poor Joint Mean of only \textbf{41.20\%}. This collapse is driven by the extreme sensitivity of its non-monotonic cosine routing equations to the high-dimensional representation shifts in the real backbone.
    \item \textbf{L3-Linear Router (Ours):} Our proposed L3-Linear router successfully generalizes to the real-world model, achieving a stable Joint Mean accuracy of \textbf{84.80\%}---outperforming QWS-Merge by a massive \textbf{+43.60\%} absolute margin.
\end{itemize}
While this empirical pilot proves that the structural insights isolated in our sandbox are highly predictive of real-world behaviors, we must temper these conclusions by acknowledging the scale and complexity limits of our pilot. Specifically, we merge only $K=3$ relatively simple visual classification tasks on a CLIP-ViT-B/16 backbone. As we scale to highly complex, diverse real-world tasks (such as ImageNet, Stanford Cars, or Flowers102) and massive LLMs, the weight-space coordinate misalignment ($\text{Error}_{alignment}$) scales significantly due to non-linear representation drift across independent fine-tuning paths. Increased task diversity and dataset complexity will inevitably introduce more pronounced representation barriers and coordinate conflicts in weight space, which could exacerbate classical routing instability if individual task trajectories diverge too wildly. Thus, while our findings demonstrate that classical routing equations scale successfully to commercial parameter manifolds under moderate task settings, future work must explore how advanced coordinate-alignment techniques (e.g., RE-Basin or Permutation Speculation) can be coupled with classical routing to mitigate scaling-induced alignment errors.
